\documentclass[12pt, a4paper]{article}

% ─── Packages ───────────────────────────────────────────────
\usepackage[margin=2.2cm]{geometry}
\usepackage[table]{xcolor}
\usepackage{titlesec}
\usepackage{booktabs}
\usepackage{tabularx}
\usepackage{colortbl}
\usepackage{array}
\usepackage{enumitem}
\usepackage{parskip}
\usepackage{microtype}
\usepackage{fancyhdr}
\usepackage{hyperref}
\usepackage[T1]{fontenc}
\usepackage[utf8]{inputenc}
\usepackage{lmodern}

% ─── Colors ─────────────────────────────────────────────────
\definecolor{accent}{HTML}{185FA5}
\definecolor{accent2}{HTML}{E85D24}
\definecolor{lightbg}{HTML}{EEF4FB}
\definecolor{midgray}{HTML}{888780}
\definecolor{darktext}{HTML}{2C2C2A}
\definecolor{rowalt}{HTML}{F5F8FC}
\definecolor{tableborder}{HTML}{D3D1C7}

% ─── Hyperref ────────────────────────────────────────────────
\hypersetup{
    colorlinks=true,
    linkcolor=accent,
    urlcolor=accent,
    pdftitle={IDSS Phase 1 - Problem Statement},
    pdfauthor={Your Team Name}
}

% ─── Section Formatting ─────────────────────────────────────
\titleformat{\section}
  {\large\bfseries\color{accent}}
  {\thesection.}{0.5em}{}
  [\vspace{-4pt}\textcolor{tableborder}{\rule{\linewidth}{0.4pt}}\vspace{2pt}]

\titleformat{\subsection}
  {\normalsize\bfseries\color{darktext}}
  {\thesubsection.}{0.5em}{}

\titlespacing{\section}{0pt}{10pt}{4pt}
\titlespacing{\subsection}{0pt}{6pt}{2pt}

% ─── Table Column Types ──────────────────────────────────────
\newcolumntype{L}[1]{>{\raggedright\arraybackslash}p{#1}}
\newcolumntype{C}[1]{>{\centering\arraybackslash}p{#1}}
\newcolumntype{Y}{>{\raggedright\arraybackslash}X}
\setlength{\extrarowheight}{1pt}
\setlength{\parskip}{3pt}
\setlength{\parindent}{0pt}

% ─── Header/Footer ───────────────────────────────────────────
\pagestyle{fancy}
\fancyhf{}
\renewcommand{\headrulewidth}{0.4pt}
\renewcommand{\headrule}{\color{tableborder}\hrule width\headwidth height\headrulewidth}
\fancyhead[L]{\small\color{midgray} IDSS Final Project --- Phase 1}
\fancyhead[R]{\small\color{midgray} Online Retail Dataset}
\fancyfoot[C]{\small\color{midgray}\thepage}

% ─── Table Header Style ──────────────────────────────────────
\newcommand{\thead}[1]{\textcolor{white}{\textbf{#1}}}

% ─── Bullet List Style ───────────────────────────────────────
\setlist[itemize]{
    leftmargin=1.4em,
  itemsep=1pt,
  topsep=2pt,
    parsep=0pt
}

% ════════════════════════════════════════════════════════════
\begin{document}
% ════════════════════════════════════════════════════════════

% ─── Title Block ─────────────────────────────────────────────
{\color{midgray}\small IDSS Final Project --- Phase 1}\\[2pt]
{\color{accent}\LARGE\bfseries Problem Statement \& Exploratory Findings}\\[4pt]
{\color{midgray}\small Online Retail Dataset (UCI ML Repository ID=352) \quad|\quad Customer Churn \& Purchase Behaviour Prediction}\\[2pt]
\textcolor{tableborder}{\rule{\linewidth}{0.5pt}}

\vspace{4pt}

% ════════════════════════════════════════════════════════════
\section{Real-World Problem}
% ════════════════════════════════════════════════════════════

A UK-based online gift retailer sells exclusively via non-store channels. Despite holding a growing
transaction history, the business currently applies identical post-purchase treatment to all
customers, leaving significant retention opportunities unexploited. The core problem is:
\textbf{the business cannot identify which customers are drifting toward churn} until it is too
late to intervene cost-effectively.

The proposed IDSS addresses this by predicting, at the customer level, the probability of
churn---defined as no purchase in the 90 days preceding the observation date---using
transactional RFM signals and engineered behavioural features.

% ════════════════════════════════════════════════════════════
\section{Stakeholders \& Decision-Making Needs}
% ════════════════════════════════════════════════════════════

\vspace{4pt}
\renewcommand{\arraystretch}{1.2}
\begin{tabularx}{\linewidth}{
  >{\columncolor{lightbg}}L{3.3cm}
  Y
  Y
}
\rowcolor{accent}
\thead{Stakeholder} & \thead{Primary Decision} & \thead{IDSS Output Needed} \\
Marketing Manager &
  Which customers to target with retention campaigns? &
  Churn probability score + risk tier \\
\rowcolor{rowalt}
Customer Success &
  Who needs proactive outreach this week? &
  Daily ranked list of at-risk customers \\
CFO / Finance &
  What is our projected customer lifetime value? &
  Segment-level projected CLV and risk-adjusted customer value \\
\end{tabularx}

% ════════════════════════════════════════════════════════════
\section{Business Value}
% ════════════════════════════════════════════════════════════

Retaining an existing customer costs 5--7$\times$ less than acquiring a new one.
The dataset spans 12 months (December 2010 -- December 2011) with \textbf{4,338 unique
customers} across \textbf{37 countries}.

\begin{itemize}
    \item Even a \textbf{10\% improvement} in retention can yield an estimated
          \textbf{\pounds 50,000--\pounds 200,000 in additional annual revenue}
          (estimated from average order value and purchase frequency).
    \item Targeted campaigns directed by churn scores reduce wasted outreach by an estimated
          \textbf{30--40\%}, improving marketing ROI.
    \item Identification of high-value at-risk customers enables differential treatment
          (personalised discounts, priority service), protecting the highest-margin segment.
\end{itemize}

% ════════════════════════════════════════════════════════════
\section{IDSS Objective \& Task Definition}
% ════════════════════════════════════════════════════════════

\textbf{Task Type:} Binary Classification

Each customer is labelled as \textbf{Churned = 1} (no purchase in the final 90 days) or
\textbf{Active = 0}. The model outputs a probability score in $[0,\,1]$ and a binary label.

\vspace{6pt}
\renewcommand{\arraystretch}{1.2}
\begin{tabularx}{\linewidth}{L{3.8cm} Y}
\rowcolor{accent}
\thead{Metric} & \thead{Justification} \\
F1-Score \emph{(primary)} &
  Balances false negatives (missed churners $\to$ lost revenue) and false positives (wasted
  marketing spend). Best suited for imbalanced binary business problems. \\
\rowcolor{rowalt}
ROC-AUC \emph{(secondary)} &
  Threshold-agnostic measure of discrimination---allows the business to set its own
  risk tolerance post-deployment. \\
Recall \emph{(tertiary)} &
  Slightly preferred: a false negative costs $\approx$\pounds 300 in lost revenue vs.\
  $\approx$\pounds 2--5 per false positive in campaign spend. \\
\end{tabularx}

\vspace{4pt}

% ════════════════════════════════════════════════════════════
\section{Dataset Overview}
% ════════════════════════════════════════════════════════════

The \textbf{Online Retail dataset} (UCI ID=352) contains 541,909 transactions from a UK-based
non-store retailer (1 December 2010 -- 9 December 2011), exposed as 6 feature columns plus
\texttt{InvoiceNo} and \texttt{StockCode} as index columns. After removing returns,
cancellations, and anonymous transactions (no CustomerID), \textbf{392,669 valid rows} across
\textbf{4,338 unique customers} remain for modelling.

\vspace{6pt}
\renewcommand{\arraystretch}{1.15}
\begin{tabularx}{\linewidth}{L{2.8cm} L{2.3cm} Y}
\rowcolor{accent}
\thead{Feature} & \thead{Type} & \thead{Notes / Quality Issues} \\
InvoiceNo      & Categorical & Index column. Prefix \texttt{C} = cancellation $\to$ excluded \\
\rowcolor{rowalt}
StockCode      & Categorical & Index column. Product ID; some codes are service charges \\
Description    & Text        & $\approx$0.3\% missing $\to$ rows dropped \\
\rowcolor{rowalt}
Quantity       & Numeric     & 10,600 rows negative (returns) $\to$ excluded \\
InvoiceDate    & Datetime    & Stored as string $\to$ converted to \texttt{datetime} \\
\rowcolor{rowalt}
UnitPrice      & Numeric     & Rows with UnitPrice $\leq 0$ $\to$ excluded \\
CustomerID     & Numeric     & $\approx$25\% missing (guest purchases) $\to$ excluded \\
\rowcolor{rowalt}
Country        & Categorical & \textbf{37} countries represented; $\approx$91\% UK \\
\end{tabularx}

% ════════════════════════════════════════════════════════════
\section{Top 3 Patterns \& Business Significance}
% ════════════════════════════════════════════════════════════

\subsection{Recency is the Dominant Churn Signal}

Recency (days since last purchase) shows the strongest correlation with the churn label
($|r| \approx 0.65$). Active customers have a mean recency of \textbf{32 days} vs.\
\textbf{214 days} for churned customers (median: 25 vs.\ 205 days).
\textit{Business implication:} A simple recency-based trigger---an automated re-engagement
email at day 60---could preemptively reach at-risk customers before they pass the churn
threshold.

\subsection{Frequency $\times$ Monetary Identifies VIP Customers}

Active customers average \textbf{121 invoices} and \textbf{\pounds 2,726} in spend vs.\
\textbf{33 invoices} and \textbf{\pounds 714} for churned customers. Frequency and monetary
spend are highly correlated ($r \approx 0.70$). \textit{Business implication:} A loyalty
programme targeting the top frequency quartile protects the highest-value segment and
maximises retention ROI.

\subsection{Geographic Variation in Churn Rate}

The dataset covers \textbf{37 countries}, with $\approx$91\% of customers based in the UK.
Several non-UK markets exhibit churn rates 10--20 percentage points above the overall mean.
\textit{Business implication:} Country-specific retention strategies (localised
communications, free-return promotions for high-churn markets) are warranted before broad
international marketing investment.

\vspace{\fill}
\textcolor{tableborder}{\rule{\linewidth}{0.4pt}}
{\small\color{midgray}
Phase 1 of 5 \quad|\quad IDSS Final Project \quad|\quad
Due: Saturday 2 May \quad|\quad Dataset: UCI Online Retail (ID=352)}

\end{document}